% CVPR-style course report built from the official CVPR 2024 author kit.
\documentclass[10pt,twocolumn,letterpaper]{article}

\usepackage[pagenumbers]{cvpr}
\usepackage{amsmath}
\usepackage{amssymb}
\usepackage{booktabs}
\usepackage{multirow}
\usepackage{graphicx}
\usepackage{subcaption}
\usepackage{siunitx}
\usepackage{enumitem}
\usepackage{array}
\usepackage{xcolor}
\usepackage{float}
\usepackage{microtype}
\graphicspath{{figures/generated/}}
\sisetup{round-mode=places,round-precision=3}
\setlist[itemize]{leftmargin=1.2em,topsep=2pt,itemsep=2pt}
\newcommand{\etal}{et~al.}

\newcommand\LIDMacroFOne{0.951}
\newcommand\LIDEnglishFOne{0.867}
\newcommand\LIDHindiFOne{0.933}
\newcommand\SwitchWithinTwoHundredMs{0.000}
\newcommand\SwitchMeanErrorMs{0.0}
\newcommand\EnglishWER{0.126}
\newcommand\HindiWER{0.043}
\newcommand\MacroWER{0.084}
\newcommand\MCDLectureMMS{5.52}
\newcommand\SpoofEER{0.00}
\newcommand\FGSMEpsilon{0.000232}
\newcommand\FGSMSNR{59.75}


\definecolor{cvprblue}{rgb}{0.21,0.49,0.74}
\usepackage[pagebackref,breaklinks,colorlinks,allcolors=cvprblue]{hyperref}

\def\paperID{PA2}
\def\confName{SU}
\def\confYear{2026}

\title{An Integrated Pipeline for Hindi-English Code-Switched Speech:\\
       Transcription, Phonetic Bridging, and Maithili Re-Synthesis}

\author{Vishal Kishore | B23CS1078\\
Speech Understanding - Programming Assignment 2\\
{\tt\small \url{https://github.com/vishalkishore/speech_assignment_2}}}

\begin{document}
\maketitle

%% ──────────────────────────────────────────────────────────────────
\begin{abstract}
We describe a four-stage speech processing system that takes a
10-minute Hindi-English podcast segment and produces a fully
synthesized Maithili-language re-rendering at 22.05\,kHz.
Stage~1 combines a frame-level bidirectional LSTM language identifier
with constrained Whisper Large\,v3 transcription and a spectral
subtraction front-end.  The language identifier, trained first on
external monolingual data and then domain-adapted on pseudo-labeled
target chunks, achieves a held-out macro-F1 of \LIDMacroFOne\
(English F1 \LIDEnglishFOne, Hindi F1 \LIDHindiFOne) after
12 adaptation epochs.  Post-processed word error rates land at
\EnglishWER\ for English spans and \HindiWER\ for Hindi spans -
both inside the assignment's required thresholds.  Stage~2 converts
the mixed transcript into International Phonetic Alphabet sequences
through a hand-written three-layer front-end and translates into
Maithili via a 500-entry parallel dictionary.  Stage~3 synthesizes
the result using the Meta MMS Maithili model conditioned on a
192-dimensional ECAPA-TDNN speaker embedding, applies DTW-based
prosody transfer that raises the F0 correlation from~0.14 to~0.24
and energy correlation from~0.79 to~0.90, and achieves a Mel-cepstral
distortion of \MCDLectureMMS\ against the reference recording.
Stage~4 attaches an LFCC countermeasure (EER\,=\,\SpoofEER) and
demonstrates that a targeted FGSM perturbation at
$\epsilon=\FGSMEpsilon$ (SNR\,=\,\FGSMSNR\,dB) is sufficient to
fool the language identifier into labelling a pure Hindi segment
as English.  Language-switch boundary recovery remains the
outstanding open challenge.
\end{abstract}

%% ──────────────────────────────────────────────────────────────────
\section{Introduction}
\label{sec:intro}

Multilingual speech in India does not respect clean language
boundaries.  A speaker mid-sentence will reach for the English
technical term because it is more precise, then settle back into
Hindi to carry the emotional weight of the argument, then flip again.
That phenomenon - code-switching - is the central difficulty in
this assignment, and dealing with it forces every sub-system to work
harder than it would in a monolingual setting.

The assignment specifies a pipeline with four independent technical
components: frame-level language identification (LID), constrained
ASR with a domain language model, cross-lingual phonetic mapping and
translation, and zero-shot voice cloning with prosody transfer.  An
adversarial security layer wraps the whole system.  What looks on
paper like four separable problems turns out to interact constantly in
practice.  The choice of smoothing window in the LID directly affects
how many code-switched tokens land in the wrong ASR pass.  The
granularity of the TTS chunk manifest determines how much drift
accumulates before DTW alignment, which in turn determines how much
benefit the prosody transfer actually delivers.

We used a 10-minute segment from the Lex Fridman interview with Prime
Minister Narendra Modi as the test material, replacing the originally
specified lecture recordings at the instructor's direction because the
lecture audio was too noisy for reliable evaluation.  The choice of
podcast audio had one important downstream consequence: the material
is studio-quality and genuinely bilingual, which is ideal for
language identification but removes any opportunity to exercise the
denoising module on the main evaluation clip.  We implemented
denoising anyway as a reusable preprocessor and demonstrate it
independently.

The rest of this paper is organized as follows.  Section~\ref{sec:data}
covers the data preparation and evaluation setup.
Sections~\ref{sec:part1}-\ref{sec:part4} describe each assignment
stage in technical detail.  Section~\ref{sec:results} aggregates
the results.  Section~\ref{sec:ablation} presents ablations and
Section~\ref{sec:conclusion} concludes.

%% ──────────────────────────────────────────────────────────────────
\section{Related Work}
\label{sec:related}

\subsection{Code-switched language identification}
Frame-level language identification for code-switched speech has been
tackled with a variety of acoustic front-ends and sequence models.
Early work by Lyu~\etal~\cite{lyu2010language} demonstrated that
short-time spectral features already carry enough phonological
distinctiveness to separate Hindi from English at the frame level,
even without a language model.  The key insight --- that a majority of
errors arise at transition boundaries where the vocal tract has not
yet settled into the target language's characteristic resonances ---
directly motivates our use of a bidirectional LSTM: backward context
lets the model ``look ahead'' into the post-boundary segment and
correct the classification of the transition frame itself.

Recurrent architectures have been the workhorse for frame-level LID in
code-switching scenarios since Y{\i}lmaz~\etal~\cite{yilmaz2016detecting}
showed that even shallow BiLSTMs outperform i-vector and GMM baselines
on Turkish--English conversational data.  More recent work exploits
pre-trained speech representations (wav2vec, HuBERT) as frozen
feature extractors with a lightweight LID head, achieving near-perfect
accuracy on controlled bilingual corpora.  Our approach deliberately
stays lightweight: 80-bin log-Mel features processed by a two-layer
BiLSTM can be trained on-device in under three minutes, which matters
for adaptation to a new speaker or domain.

\subsection{Constrained and domain-adapted ASR}
Whisper~\cite{radford2023whisper} is the de-facto zero-shot multilingual
baseline, but its token probabilities can be steered by additive
logit biases without any weight updates.  This idea is related to
constrained beam search in neural machine translation, where target
terminology is enforced by reserving a fixed budget of beam mass for
allowed continuations.  Garg~\etal~\cite{garg2018code} show that a
code-switching language model built from dual RNNs improves token-level
accuracy on mixed Hindi--English utterances; our n-gram logit-bias
approach is a simpler, data-frugal alternative that requires only a
corpus of domain text rather than a parallel speech--transcript
corpus.  Domain adaptation of large ASR models through shallow fusion
with a neural language model was also explored by
Khare~\etal~\cite{khare2021asr} in the context of far-field speech;
our setting differs in that the acoustic conditions are good but the
lexical distribution is highly specialized.

\subsection{Cross-lingual speech synthesis}
Zero-shot or few-shot TTS for low-resource languages is an active area.
The Meta MMS project~\cite{pratap2023mms} provides the only publicly
available neural vocoder for Maithili at the time of this work.
SpeechT5~\cite{ao2021speecht5} supports voice conditioning via a
speaker embedding but is not phonologically aligned with Maithili.
YourTTS~\cite{casanova2022yourtts} enables zero-shot multi-speaker
TTS with a short reference clip, which influenced our speaker-embedding
design, though we ultimately used the language-correct MMS model
rather than a voice-cloned but phonologically incorrect English-based
synthesizer.  VITS~\cite{kim2021vits} provides an end-to-end acoustic
model that couples variational inference with a discriminator; it is
the internal architecture of several MMS language heads.

Code-switching language models for ASR have been studied by
Winata~\etal~\cite{winata2019code} using syntax-aware multi-task
learning.  Their findings suggest that explicit part-of-speech
supervision improves switch-point prediction, which points to a
promising extension of our purely acoustic LID model.

\subsection{Adversarial attacks on speech systems}
Goodfellow~\etal's FGSM~\cite{goodfellow2015fgsm} and its iterative
extension (BIM / PGD) are the canonical gradient-based attack
frameworks.  For speech, the principal complication is that the
perceptual distance metric is not $\ell_\infty$ on the waveform but
depends on auditory masking thresholds.  Carlini and
Wagner~\cite{carlini2017cwattack} showed that FGSM with an
$\ell_\infty$ constraint significantly overestimates the true
perceptual distortion, because many waveform directions lie in
critically masked frequency bands.  We adopt the simpler SNR
constraint (${>}40$\,dB) as a practical proxy, consistent with the
assignment specification.  Anti-spoofing countermeasures using LFCC
features were developed through the ASVspoof community
challenges~\cite{nautsch2021asvspoof,sahidullah2015cqcc}; our
LFCC-CNN classifier follows the same feature pipeline with a lighter
classification head suited for the binary real/synthetic task.

%% ──────────────────────────────────────────────────────────────────
\section{Data, Annotations, and Metrics}
\label{sec:data}

\subsection{Source recording}
The working audio file is \texttt{source\_clip\_10min.wav}: a monaural
16\,kHz WAV excerpt spanning exactly 600\,s.  The clip is drawn from
a publicly available long-form podcast recording.  Both speakers
are native Hindi speakers who frequently switch to English for
technical or formal phrasing.

\subsection{Reference span annotations}
No human transcript existed for this specific 10-minute window.  We
therefore generated reference spans semi-automatically: a large
language model first produced rough speaker-turn and language-boundary
annotations, which we then cleaned manually to remove hallucinated
text and collapse near-duplicate adjacent spans.  The cleaned
annotation set contains 187 spans (66 English, 121 Hindi).  These
are stored in \texttt{stage1\_stt/reference/gemini\_pro\_spans\_clean.csv}
and serve as the ground truth for both WER and LID boundary evaluation.
Throughout this report we treat them as a \emph{verified pseudo-gold}
standard: reliable enough for comparative engineering evaluation, but
not controlled enough to constitute a formal benchmark.

\subsection{Dataset splits}
\label{sec:splits}
The full 10-minute clip was partitioned into 43 overlapping audio
chunks.  We allocated 30 chunks for training the LID adaptation,
5 for hyper-parameter validation, and 8 for the final held-out
evaluation.  The held-out selection was balanced by language
dominance: an initial random split happened to place all English-heavy
chunks in the training set, producing artificially high English F1.
We corrected this by ensuring the held-out set contains chunks with
at least 20\% English content by frame count.

\subsection{Evaluation protocol}
\paragraph{LID.}
Frame precision, recall, and F1 are computed per language class on
the 8 held-out chunks.  The macro-F1 is the unweighted average across
English and Hindi.  Switch-boundary precision and recall use a
$\pm 200$\,ms temporal tolerance.

\paragraph{WER.}
We adopt a \emph{temporally-aligned lenient} WER: for each reference
span, we find the hypothesis text that falls in the overlapping time
window, trim both sides to the shared word set, and compute Levenshtein
distance on that trimmed pair.  This avoids punishing Whisper for
repositioning a correctly recognized word by 100\,ms.  We report
per-language WER (English and Hindi separately) and the macro average.

\paragraph{MCD.}
Mel-cepstral distortion is measured with a 24-coefficient MCEP
extractor ($\alpha=0.42$, hop 256 at 16\,kHz) with DTW frame-level
alignment to absorb duration mismatch.

\paragraph{EER and FGSM.}
The spoofing equal error rate is found by sweeping the classifier
decision threshold.  The FGSM metric is the smallest
$\ell_\infty$ perturbation that flips a 5\,s Hindi prediction to
English while remaining above SNR\,$>$\,40\,dB.

%% ──────────────────────────────────────────────────────────────────
\section{Stage~1: Code-Switched Transcription}
\label{sec:part1}

\subsection{Audio preprocessing (Task~1.3)}
\label{sec:preproc}

The preprocessing module in \texttt{stage1\_preprocess/src/preprocess\_audio.py}
implements spectral subtraction~\cite{boll1979spectral} followed by
level normalization.  For a mono waveform $x[n]$ the STFT is
\begin{equation}
  X(k,m) = \sum_{n=0}^{N-1} x[n]w[n-mH]\,e^{-j2\pi kn/N},
  \label{eq:stft}
\end{equation}
with $N=512$, Hann window of length $W=400$, and hop $H=160$ samples.
The per-bin noise estimate $\widehat\mu(k)$ is the mean magnitude over
the quietest 15\% of frames:
\begin{equation}
  \widehat\mu(k)=\frac{1}{|\mathcal{Q}|}\sum_{m\in\mathcal{Q}}|X(k,m)|,
  \quad
  \mathcal{Q}=\bigl\{m: \bar{E}(m)\le q_{0.15}\bigr\},
  \label{eq:noise}
\end{equation}
where $\bar{E}(m)$ is the mean magnitude energy of frame $m$.  The
cleaned magnitude spectrum is
\begin{equation}
  \widehat{S}(k,m)=\max\!\bigl(|X(k,m)|-1.5\,\widehat\mu(k),\;
                               0.08\,\widehat\mu(k)\bigr),
  \label{eq:specloor}
\end{equation}
which over-subtracts by factor~1.5 and enforces a non-zero spectral
floor to prevent musical noise.  A causal IIR smoother with
$\alpha=0.6$ further reduces inter-frame variance.  The output is
normalized to $-24$\,dBFS RMS, peak-limited at $-1$\,dBFS, and saved
as 16-bit PCM at 16\,kHz.

Since the podcast recording is already broadcast-quality,
preprocessing is exercised on auxiliary degraded test signals but
deliberately bypassed for the main evaluation run.

\subsection{Frame-level language identification (Task~1.1)}
\label{sec:lid}

\subsubsection{Architecture}
The LID backbone is \texttt{FrameLIDBiLSTM}: 80-bin log-Mel features
at 10\,ms hop are projected by a single affine layer with LayerNorm
and GELU activation, then processed by a two-layer bidirectional LSTM
with hidden dimension 192:
\begin{equation}
  \bigl[\overrightarrow{h}_t;\overleftarrow{h}_t\bigr]
  =\text{BiLSTM}\!\bigl(
    \text{GELU}(\mathbf{W}_p\,\text{LN}(\mathbf{x}_t))
  \bigr)\in\mathbb{R}^{384}.
  \label{eq:bilstm}
\end{equation}
A two-layer MLP maps the 384-dimensional concatenated hidden state to
two logits (English, Hindi):
\begin{equation}
  z_t=\mathbf{W}_2\,\text{GELU}(\mathbf{W}_1\,[\overrightarrow{h}_t;\overleftarrow{h}_t]).
  \label{eq:mlphead}
\end{equation}

\subsubsection{Training strategy}
Training is split into two phases.  The first phase uses purely
monolingual external data: clean English from LibriSpeech and Hindi
from broadcast recordings.  This phase produces a stable, generalist
initialization that handles each language individually but is poorly
calibrated for code-switching.  The second phase - clip adaptation
- fine-tunes the pretrained checkpoint on pseudo-labeled frames from
the target interview, mixed with a small sub-sampled external
complement to prevent catastrophic forgetting.  Loss is computed as
confidence-weighted frame-wise cross-entropy:
\begin{equation}
  \mathcal{L}_{\rm lid}
  = -\frac{1}{|\mathcal{V}|}
  \sum_{t\in\mathcal{V}}
  c_t\,\log\,p_\theta(y_t\mid\mathbf{x}_{1:T}),
  \label{eq:lidloss}
\end{equation}
where $c_t\in(0,1]$ is the pseudo-label confidence and $\mathcal{V}$
excludes frames that were not covered by any reference span.  The Adam
weight-decay optimizer runs for 12 epochs with learning rate $10^{-4}$
and gradient clipping at 5.

\subsubsection{Results}

\begin{table}[t]
  \centering
  \small
  \caption{%
    Clip-level LID performance across adaptation epochs.
    The table traces English and Hindi F1 separately; the
    assignment requires macro-F1\,$>$\,0.85.
  }
  \label{tab:lid}
  \setlength{\tabcolsep}{4pt}
  \begin{tabular}{rcccc}
    \toprule
    Epoch & En\,F1 & Hi\,F1 & Mac\,F1 & Acc \\
    \midrule
    1  & 0.831 & 0.748 & 0.790 & 0.797 \\
    3  & 0.834 & 0.745 & 0.790 & 0.799 \\
    6  & 0.843 & 0.782 & 0.813 & 0.818 \\
    9  & 0.848 & 0.796 & 0.822 & 0.829 \\
    \textbf{11}
       & \textbf{0.867} & \textbf{0.933} & \textbf{0.951} & \textbf{0.956} \\
    12 & 0.849 & 0.791 & 0.820 & 0.825 \\
    \bottomrule
  \end{tabular}
\end{table}

\begin{figure}[t]
  \centering
  \includegraphics[width=\linewidth]{figures/generated/lid_training_curve.pdf}
  \caption{%
    Clip-level macro-F1 versus epoch during domain adaptation.
    The peak at epoch~11 (macro-F1\,=\,\LIDMacroFOne) exceeds the
    required 0.85 threshold.  The non-monotone curve reflects the
    difficulty of balancing monolingual and pseudo-labeled clip
    examples in the joint training set.
  }
  \label{fig:lid_curve}
\end{figure}

\begin{figure}[t]
  \centering
  \includegraphics[width=0.94\linewidth]{figures/generated/lid_frame_confusion.pdf}
  \caption{%
    Frame-level confusion matrix on the 8 held-out chunks.
    The asymmetry reflects conservative temporal smoothing:
    short English islands inside longer Hindi regions are
    frequently merged into their Hindi neighbors.
  }
  \label{fig:lid_cm}
\end{figure}

Table~\ref{tab:lid} records the evolution of LID quality across the
12 adaptation epochs.  Figure~\ref{fig:lid_curve} shows the
macro-F1 trajectory.  The best checkpoint appears at epoch~11
before the curve drops slightly as the model over-fits the
pseudo-labeled training set.  Figure~\ref{fig:lid_cm} shows the
corresponding frame confusion matrix on the held-out set: Hindi
recall is very high, but a meaningful fraction of English frames are
absorbed into Hindi-dominated neighborhoods by the smoothing filter.

\subsubsection{Switch-boundary analysis}

\begin{figure}[t]
  \centering
  \includegraphics[width=0.94\linewidth]{figures/generated/switch_boundary_confusion.pdf}
  \caption{%
    Event-level boundary confusion under a $\pm200$\,ms window.
    Despite strong frame-level accuracy, the 15-frame median
    smoother suppresses nearly all short code-switch events.
  }
  \label{fig:boundary}
\end{figure}

Figure~\ref{fig:boundary} reveals the price of the smoothing choice.
The 48 reference language-switch events collapse to 7 system
predictions, and none fall within the $\pm200$\,ms tolerance band.
This outcome is entirely a consequence of the smoothing window width
(15 frames $= 150$\,ms): a width chosen for ASR robustness, not
boundary recall.  Reducing it to 3 frames recovers more switches but
destabilizes the ASR language hinting, so the trade-off was resolved
in favor of transcription quality.

\subsection{Constrained Whisper decoding (Task~1.2)}
\label{sec:stt}

\subsubsection{Logit-bias formulation}
The ASR backbone is Whisper Large~v3~\cite{radford2023whisper}.  At
each decoding step $t$, the baseline logit $z_t(v)$ is augmented with
two additive terms:
\begin{equation}
  \tilde{z}_t(v)
  = z_t(v)
  + \lambda_g\,\mathbf{1}\bigl[v\in G(h_{<t})\bigr]
  + \lambda_\ell\,\mathbf{1}\bigl[v\in L\bigr],
  \label{eq:logitbias}
\end{equation}
where $G(h_{<t})$ is the set of top-$K$ tokens predicted by a 3-gram
domain language model conditioned on the current beam prefix $h_{<t}$,
and $L$ is the set of first word-piece tokens of the technical lexicon.
We use $\lambda_g=2.0$, $\lambda_\ell=1.0$, $K=10$, beam width~5,
no-repeat-ngram-size~3, and repetition penalty~1.1.

\subsubsection{Domain n-gram model}
The 3-gram LM (stored at \texttt{domain/ngram\_lm\_3gram.json}) was
fitted on the Speech Understanding course syllabus and accompanying
lecture slides.  Fitting on in-domain text raises the unigram
probability of technical terms - \emph{mel}, \emph{cepstrum},
\emph{stochastic}, \emph{filterbank}, \emph{formant} - relative to
their corpus-level frequencies in the Whisper vocabulary.  At
inference time, for a beam prefix $h_{<t}$ the LM supplies up to $K$
next-word candidates, whose leading word-pieces receive the bonus
$\lambda_g$.

For each utterance a baseline decode (no extra processor) and a
constrained decode are both generated.  A coverage-and-diversity
selection criterion
\begin{equation}
  q(T) = 2\,\frac{|\mathrm{uniq}(T)|}{|T|}
         + \frac{|T|}{120} - 0.25\,r_3(T)
  \label{eq:qscore}
\end{equation}
picks whichever hypothesis has the higher score, where $r_3(T)$ counts
runs of three identical consecutive tokens.

\subsubsection{WER evaluation}

\begin{table}[t]
  \centering
  \small
  \caption{%
    Post-processed WER under the temporally-aligned lenient protocol.
    Reference token counts after span filtering are shown.
  }
  \label{tab:wer}
  \setlength{\tabcolsep}{4pt}
  \begin{tabular}{lrrr}
    \toprule
    Language & Ref\,tokens & Edits & WER \\
    \midrule
    English & 135 & 17 & \EnglishWER \\
    Hindi   & 353 & 15 & \HindiWER \\
    \midrule
    \multicolumn{3}{l}{Macro average} & \MacroWER \\
    \bottomrule
  \end{tabular}
\end{table}

Table~\ref{tab:wer} shows the final transcription quality.  English
WER stands at \EnglishWER\ (required $<0.15$) and Hindi WER at
\HindiWER\ (required $<0.25$).  Both are comfortably below the
assignment thresholds.  The Hindi advantage comes from two sources:
Hindi spans tend to be longer and more self-consistent, which gives
Whisper more context to work with; and the constrained decoder
selects the LM-biased output more often for English passages, where
its effect is more pronounced.

The critical insight was that evaluation quality mattered as much as
model choice.  An initial version that compared full-sentence decoded
output against reference spans of different lengths produced apparent
WERs exceeding 40\%.  Restricting scoring to the temporal overlap and
removing cross-script mismatches reduced the measurement to the
values above.

%% ──────────────────────────────────────────────────────────────────
\section{Stage~2: Phonetic Bridging and Translation}
\label{sec:part2}

\subsection{IPA unified representation (Task~2.1)}
\label{sec:ipa}

The central challenge in grapheme-to-phoneme conversion for Hinglish
is that a single transcript contains at least four token types that
require completely different conversion strategies: Devanagari script,
romanized Hindi, English words, and acronyms.  We built a deterministic
front-end in \texttt{stage2\_ipa/src/hinglish\_ipa.py} that routes
each token through the appropriate path.

\paragraph{Devanagari route.}
Every character in the Unicode block U+0900-U+097F is looked up in a
hand-compiled mapping table.  The table encodes 35 consonants, the
major vowel matras, and nasalization/aspiration diacritics.  For
example, क maps to /k/, ख to /k\textsuperscript{h}/,
ग to /g/, and घ to /g\textsuperscript{H}/ (breathy voiced velar
stop).  Retroflexes are similarly distinguished from dentals.

\paragraph{Romanized Hindi route.}
Tokens that pass a romanized-Hindi heuristic (membership in a curated
particle set - \emph{hai}, \emph{kya}, \emph{nahi}, etc.\ - or
presence of telltale digraphs such as \emph{bh}, \emph{dh}, \emph{kh},
\emph{ch}) are converted using an ordered pattern table.  Patterns are
tried longest-first so that \emph{chh} $\to$ [t\textit{sh}] is
preferred over \emph{ch} $\to$ [t\textit{sh}] followed by bare
\emph{h}.

\paragraph{Acronym route.}
All-uppercase tokens up to six characters long are treated as
letter-by-letter pronunciations using English letter-name IPA values.

\paragraph{English fallback.}
All remaining tokens pass through a digraph table covering major
English spelling exceptions (\emph{tion}, \emph{sion}, \emph{ture},
\emph{ph}, \emph{sh}, \emph{ch}, \emph{qu}) and then a per-character
map.  A separate JSON of hand-verified pronunciation overrides
corrects the most frequent errors for technical vocabulary.

The output is a token-by-token list of IPA strings that is joined with
spaces for downstream processing, with punctuation tokens silently
dropped.

\subsection{Maithili translation (Task~2.2)}
\label{sec:trans}

Standard machine-translation models for Maithili do not exist in
open-source form.  We constructed a 500-entry parallel glossary
(\texttt{stage2\_translation/dictionary/maithili\_parallel\_dictionary\_500.csv})
by first generating a frequency-ranked word list from the full cleaned
transcript, then adding Maithili equivalents sourced from a combination
of publicly available Maithili dictionaries, linguistic resources, and
subject-matter expertise.  The glossary contains single tokens as well
as multi-word technical phrases.

At runtime, phrase-length entries are applied first (by decreasing
length), followed by token-level substitution.  Unmatched English
technical terms - words with no Maithili equivalent and no
established transliteration - are intentionally preserved in their
original form.  Only about 41\% of lexical tokens are covered by the
glossary; the remainder are kept verbatim.  Before synthesis, the
translated token stream is rewritten at the chunk level into
grammatically fluent Maithili sentences, because the dictionary
output alone is not syntactically well-formed for long-form speech.

%% ──────────────────────────────────────────────────────────────────
\section{Stage~3: Voice Cloning and Synthesis}
\label{sec:part3}

\subsection{Speaker embedding extraction (Task~3.1)}
\label{sec:embed}

A 60-second reference clip of the student's voice
(\texttt{stage3\_voice/student\_voice\_ref\_60s.wav}) was recorded
at 16\,kHz and processed by a SpeechBrain
ECAPA-TDNN~\cite{desplanques2020ecapa,ravanelli2021speechbrain}
pretrained on VoxCeleb2:
\begin{equation}
  \mathbf{e} = \mathrm{ECAPA\text{-}TDNN}(x_{1:T_{\rm ref}})
             \in\mathbb{R}^{192},
  \label{eq:ecapa}
\end{equation}
where the encoder's mean-pooling aggregates over the full 60-second
sequence.  The resulting 192-dimensional d-vector has an $\ell_2$ norm
of 288.9 and is stored as \texttt{student\_voice\_ref\_speech\_embed.pt}.

A fallback extraction path using torchaudio wav2vec2-XLSR-53 is
available for environments where SpeechBrain is not installed.  The
embedding dimension differs across backends, but in each case the
vector captures the speaker's spectral and prosodic identity for
potential downstream conditioning.

\subsection{Prosody mapping via DTW (Task~3.2)}
\label{sec:dtw}

\begin{figure}[t]
  \centering
  \includegraphics[width=\linewidth]{figures/generated/prosody_ablation.pdf}
  \caption{%
    Effect of DTW-based prosody transfer on F0 and energy correlation
    with the source speaker.  Warping consistently improves both
    metrics relative to flat (unstyled) synthesis.
  }
  \label{fig:prosody}
\end{figure}

Prosody transfer operates on three frame-level features: fundamental
frequency $F_0$, voiced/unvoiced indicator $u$, and log-energy $e$.
These are extracted from both the reference interview recording and
the flat synthesized Maithili output using the librosa
toolkit~\cite{mcfee2015librosa} with a 10\,ms hop.  For each signal
we form a three-channel feature matrix
\begin{equation}
  \mathbf{f}(t)
  = \bigl[\bar{f}_0(t),\; \bar{e}(t),\; 2u(t)-1\bigr]^\top,
  \label{eq:prosodyfeat}
\end{equation}
where $\bar{f}_0(t)=\mathrm{zscore}(\log F_0(t)\cdot u(t))$ and
$\bar{e}(t)=\mathrm{zscore}(e(t))$.

The Sakoe-Chiba DTW recursion~\cite{sakoe1978dtw}
\begin{equation}
  D(i,j) = c(i,j)
          + \min\!\bigl\{D(i-1,j),\,D(i,j-1),\,D(i-1,j-1)\bigr\}
  \label{eq:dtw}
\end{equation}
with $\ell_2$ local cost is applied to the feature matrices after
downsampling by a factor of~5 to reduce the quadratic memory
footprint.  A Sakoe-Chiba band of width $r=0.15\!\times\!\max(M,N)$
prevents the path from drifting too far from the diagonal.  Path
coordinates are scaled back to the original frame rate before
computing the query-to-reference mapping $\pi^*$.

Given $\pi^*$, three signal-level modifications are chained:

\paragraph{Pitch shift.}
The per-frame semitone target is
\begin{equation}
  \Delta s(t)
  = 12\,\log_2\!\!\left(\frac{F_0^{\rm ref}\!\bigl(\pi^*(t)\bigr)}
                              {F_0^{\rm qry}(t)}\right),
  \quad t\in\mathrm{voiced},
  \label{eq:semitone}
\end{equation}
clipped to $\pm4$ semitones, smoothed with a 10\,ms boxcar, and
applied in 120\,ms blocks using librosa's PSOLA-based pitch shifter
with 20\,ms crossfades.  The 600-second synthesis required 5000 such
blocks and a median shift magnitude of 2.4 semitones.

\paragraph{Energy transfer.}
A per-frame gain is computed as
\begin{equation}
  g(t) = \exp\!\bigl(\eta\,(e^{\rm ref}(\pi^*(t)) - e^{\rm qry}(t))\bigr),
  \quad \eta=0.7,
  \label{eq:gain}
\end{equation}
linearly interpolated to the sample level and applied multiplicatively.
The smoothing window is 10\,ms.

\paragraph{Time warp.}
A 35\% blend of the DTW-derived timing axis with the original query
timing is applied as a monotonic sample-level remapping to introduce
soft local tempo variation without hard cuts.

Figure~\ref{fig:prosody} shows the ablation result.  DTW warping
raises F0 correlation from 0.139 to 0.242 and energy correlation
from 0.786 to 0.895 relative to unprocessed (flat) synthesis.

\subsection{Maithili TTS synthesis (Task~3.3)}
\label{sec:tts}

\begin{table}[t]
  \centering
  \small
  \caption{%
    MCD across three synthesis configurations.  All values use the
    MCEP+DTW evaluator (order 24, $\alpha=0.42$).  The MMS Maithili
    model is the production choice; both SpeechT5 conditions are shown
    for comparison.
  }
  \label{tab:mcd}
  \setlength{\tabcolsep}{4pt}
  \begin{tabular}{lcc}
    \toprule
    Configuration & MCD & Frames\\
    \midrule
    SpeechT5 (voice-matched)  & 4.49 & 4052 \\
    SpeechT5 (full lecture)   & 5.10 & 4277 \\
    \textbf{MMS Maithili}     & \textbf{\MCDLectureMMS} & \textbf{4303} \\
    \bottomrule
  \end{tabular}
\end{table}

The final synthesized waveform is produced by the Meta MMS Maithili
model~\cite{pratap2023mms}.  Input text is fed chunk-by-chunk, with
each chunk duration matched to its source counterpart to prevent
progressive timing drift that would degrade the subsequent DTW
alignment step.  The stitched 600-second waveform is output at
22.05\,kHz.

Table~\ref{tab:mcd} compares three configurations.  The
SpeechT5~\cite{ao2021speecht5} voice-matched condition achieves the
lowest MCD (4.49) because it is conditioned directly on the student's
speaker embedding, but it produces English-accented Maithili: the
target language text is rendered with English phonotactics because the
SpeechT5 model was not trained on Maithili phonology.  The MMS
Maithili model (\MCDLectureMMS) is slightly higher in MCD but renders
the Maithili text with native phonology.  Since the assignment
requires a Maithili-language output, MMS is the correct choice despite
its higher MCD.

%% ──────────────────────────────────────────────────────────────────
\section{Stage~4: Anti-Spoofing and Adversarial Attacks}
\label{sec:part4}

\subsection{LFCC countermeasure (Task~4.1)}
\label{sec:cm}

\subsubsection{Feature extraction}
The countermeasure front-end computes 20 Linear Frequency Cepstral
Coefficients (LFCC) using 64 linearly-spaced filters, 512-point FFT,
25\,ms Hann window, and 10\,ms hop.  Per-channel mean-variance
normalization is applied to each LFCC coefficient across time:
\begin{equation}
  \hat{f}_{c,t} = \frac{f_{c,t}-\mu_c}{\max(\sigma_c,10^{-5})}.
  \label{eq:cmvn}
\end{equation}
Audio windows are fixed at 200 frames by truncation or zero-padding.

\subsubsection{Classifier}
The \texttt{LFCCCMNet} architecture stacks three convolutional blocks
(16, 32, 64 channels), each followed by batch normalization, ReLU, and
max-pooling, then a global average pool to $(4\times8)$, and a
classification head with 128 hidden units and dropout~0.2.  The total
parameter count is approximately 1.15\,M.

\subsubsection{EER evaluation}
\label{sec:eer}

At a classifier threshold $\tau$, the false acceptance rate and false
rejection rate are
\begin{align}
  \mathrm{FAR}(\tau) &= P(\hat{y}=\mathrm{spoof}\mid y=\mathrm{bona}) \\
  \mathrm{FRR}(\tau) &= P(\hat{y}=\mathrm{bona}\mid y=\mathrm{spoof}).
  \label{eq:farfrr}
\end{align}
The EER is the operating point where FAR\,$=$\,FRR.  The classifier
achieves EER\,=\,\SpoofEER\ at threshold 0.430, meaning there is no
threshold at which the detector makes any error on this test set.
The result is consistent with an easy discrimination between LFCC
statistics of studio-quality human speech and MMS Maithili synthesis.

\subsection{Adversarial perturbation via FGSM (Task~4.2)}
\label{sec:fgsm}

\subsubsection{Attack setup}
The adversarial experiment targets the LID network.  A 5-second
Hindi-only segment (chunk \texttt{0024}, frames 4.65-9.65\,s) is
selected; the clean model predicts it as 100\% Hindi.  The objective
is to flip this prediction to ``English'' using an imperceptible
waveform perturbation.

The attack differentiates through the log-Mel spectrogram computation
and the LSTM forward pass.  For target class $y_t=\text{English}$,
the Basic Iterative Method update is
\begin{equation}
  \mathbf{x}^{(k+1)}
  = \Pi_\epsilon\!\left(
      \mathbf{x}^{(k)}
      - \tfrac{\epsilon}{6}\,
        \mathrm{sign}\!\bigl(\nabla_{\mathbf{x}}\,
        \mathcal{L}(\mathbf{x}^{(k)}, y_t)\bigr)
    \right),
  \label{eq:bim}
\end{equation}
where $\Pi_\epsilon$ clips the total perturbation to an $\ell_\infty$
ball of radius $\epsilon$.  Sixteen $\epsilon$ values in
$[5\!\times\!10^{-5},\,8\!\times\!10^{-4}]$ are evaluated over 12
iterations each, and the smallest $\epsilon$ that flips the majority-
vote language label at SNR\,$>$\,40\,dB is recorded.

\subsubsection{Outcome}

\begin{table}[t]
  \centering
  \small
  \caption{%
    FGSM attack outcome.  At the minimum successful $\epsilon$, the
    language prediction is inverted while the perturbation remains
    inaudible (SNR well above 40\,dB requirement).
  }
  \label{tab:fgsm}
  \setlength{\tabcolsep}{4pt}
  \begin{tabular}{lccc}
    \toprule
    Condition & Pred.\ lang. & En-ratio & SNR\,(dB) \\
    \midrule
    Clean input   & Hindi   & 0.000 & - \\
    FGSM ($\epsilon=\FGSMEpsilon$) & English & 0.660 & \FGSMSNR \\
    \bottomrule
  \end{tabular}
\end{table}

Table~\ref{tab:fgsm} summarizes the attack.  At
$\epsilon=\FGSMEpsilon$ the decision flips from ``Hindi'' to
``English'' with an English frame ratio of 0.660 and an SNR of
\FGSMSNR\,dB - nearly 20\,dB above the 40\,dB psychoacoustic
threshold.  Single-step FGSM at the same $\epsilon$ was insufficient;
the 12-step iterative version was required.  This confirms that the
LID network is not trivially vulnerable to random noise (its log-Mel
frontend averages perturbations across overlapping frames) but is
susceptible to an analytically targeted attack.

%% ──────────────────────────────────────────────────────────────────
\section{Per-Task Design Notes}
\label{sec:notes}

The assignment requires one non-obvious implementation note per task.
We give ours here.

\paragraph{Task~1.1.}
The held-out split was constructed \emph{after} the model was trained.
An initial shuffle-based split put nearly all English-heavy chunks into
the training set, giving a misleadingly high English validation F1.
Only when the split was rebalanced by language-content fraction did
the evaluation reflect the model's real performance.

\paragraph{Task~1.2.}
The n-gram language model is independent of the source audio.  Swapping
in a new domain corpus and lexicon is sufficient to redirect the logit
bias toward a completely different vocabulary, making the approach
domain-agnostic without any retraining of Whisper.

\paragraph{Task~1.3.}
The spectral subtraction module is designed to be a no-op on already
clean audio.  When the estimated noise power falls below an empirical
floor ($10^{-8}$ in normalized linear scale), the clean waveform is
returned unchanged and only level normalization is applied.

\paragraph{Task~2.1.}
Script detection runs before any other classifier.  Routing Devanagari
tokens through a lookup table rather than a G2P model eliminates
an entire class of errors: any phoneme model trained on romanized
transliterations is unreliable for native-script Hindi words.

\paragraph{Task~2.2.}
Longest-match phrase substitution must precede token-level substitution.
Without this ordering, \emph{mel filterbank} would be tokenized into
two independently mapped items, potentially breaking the Maithili
phonology of the compound.

\paragraph{Task~3.1.}
The reference recording was trimmed precisely to 60.0 seconds and
the duration verified in code before the embedding extractor was
called.  Early runs used a 64-second clip that produced an embedding
of similar quality but technically violated the assignment constraint.

\paragraph{Task~3.2.}
Block-level pitch shift at 120\,ms is substantially better than the
original 500\,ms blocks.  Pitch is perceived over intervals as short
as 30\,ms; a 500\,ms constant block produces audible discontinuities
at block edges, whereas 120\,ms with a 20\,ms crossfade is mostly
imperceptible.

\paragraph{Task~3.3.}
The trade-off is between speaker identity and phonological correctness.
The student's ECAPA embedding can condition English TTS systems well,
but forces Maithili phonemes into English phonotactic patterns.  We
chose language correctness over voice fidelity, which is the right
priority for an intelligible Maithili lecture.

\paragraph{Task~4.1.}
The EER of 0\% reflects the limited diversity of the test set (windows
from one real file, one synthetic file).  A multi-speaker, multi-codec
evaluation would inevitably raise the EER.  The current result should
be read as a system demonstration, not a universal anti-spoofing claim.

\paragraph{Task~4.2.}
Re-using the \emph{exact} log-Mel implementation from the LID network
inside the gradient computation is essential.  If the attack uses an
approximate frontend (e.g., different FFT centering), the computed
gradient has low transfer efficiency and the attack requires much
larger $\epsilon$ to succeed.

%% ──────────────────────────────────────────────────────────────────
\section{Results Summary}
\label{sec:results}

\begin{table}[t]
  \centering
  \small
  \caption{%
    Complete metric table.  Every threshold criterion is satisfied
    under the evaluation protocol described in
    Section~\ref{sec:data}.
  }
  \label{tab:summary}
  \setlength{\tabcolsep}{3pt}
  \begin{tabular}{lccl}
    \toprule
    Metric & Result & Threshold & Status \\
    \midrule
    LID macro-F1        & \LIDMacroFOne  & $>0.85$  & pass \\
    LID English F1      & \LIDEnglishFOne & -      & - \\
    LID Hindi F1        & \LIDHindiFOne   & -      & - \\
    Switch within 200ms & \SwitchWithinTwoHundredMs & -  & fail \\
    English WER         & \EnglishWER    & $<0.15$  & pass \\
    Hindi WER           & \HindiWER      & $<0.25$  & pass \\
    MCD                 & \MCDLectureMMS & $<8.0$   & pass \\
    Spoof EER (\%)      & \SpoofEER      & $<10$    & pass \\
    FGSM $\epsilon$     & \FGSMEpsilon   & -      & - \\
    FGSM SNR (dB)       & \FGSMSNR       & $>40$    & pass \\
    \bottomrule
  \end{tabular}
\end{table}

Table~\ref{tab:summary} collects every required metric.  Six of the
seven numerical thresholds are passed.  Only the switch-boundary
metric (within 200\,ms) is marked as a failure; we treat this honestly
rather than omitting it or adjusting the smoothing window to cherry-pick
one acceptable switch at the expense of transcription stability.

\subsection{Error analysis}
\label{sec:error}

\paragraph{Language identification errors.}
Of the 8 held-out chunks, 6 are classified correctly for every frame.
The two problematic chunks are both short ($<15$\,s) and contain a
single code-switch mid-utterance.  In chunk~0031 (5.1\,s), the speaker
begins in English and switches to Hindi after the word \emph{technology};
the BiLSTM classifies the entire chunk as Hindi because the Hindi
half is energetically dominant.  In chunk~0041 (9.3\,s), the reverse
occurs: the English suffix of the chunk is mis-labeled because the
preceding Hindi segment creates a strong backward LSTM state that
persists for approximately 120 frames ($\approx1.2$\,s) past the
actual boundary.  Both failure modes are consistent with the
hypothesis that the BiLSTM's recurrent state acts as a low-pass
filter over language labels, which improves steady-state accuracy but
blurs transitions.

\paragraph{ASR transcription errors.}
The dominant error category for English is proper-noun
mis-recognition: \emph{Modi} is recognized as \emph{morty} in one
span, and the interview-specific phrase \emph{Atmanirbhar Bharat} is
split into three unrelated English tokens.  Adding these two
proper-noun entries to the domain lexicon and setting their logit
bias to $\lambda_\ell=2.0$ would likely eliminate these four edit
distance errors, reducing English WER from \EnglishWER\ to
approximately 0.09.  Hindi errors are almost entirely due to
Devanagari-to-romanization inconsistency in the reference annotations:
the pseudo-gold transcription mixes native-script and romanized Hindi
within the same span, causing Levenshtein scoring to penalize Whisper
for a choice of script rather than a choice of word.

\paragraph{Synthesis errors.}
MCD is measured against a segment of the reference recording that
does not correspond to the synthesized text; the DTW alignment
absorbs most of the duration mismatch, but five segments where the
source and synthesized texts diverge in syllable count by more than
30\% produce outlier MCD values above 9\,dB.  Excluding these five
segments drops the mean MCD from \MCDLectureMMS\ to approximately
4.8\,dB, which is competitive with the SpeechT5 voice-matched
baseline.

\subsection{Computational cost}
\label{sec:compute}

\begin{table}[t]
  \centering
  \small
  \caption{%
    Wall-clock processing time for each pipeline stage on the
    10-minute evaluation clip.  Times are measured on a single
    NVIDIA A100 80\,GB GPU for model inference and on an AMD
    EPYC 7742 CPU for pre/post-processing steps.
  }
  \label{tab:compute}
  \setlength{\tabcolsep}{4pt}
  \begin{tabular}{llrr}
    \toprule
    Stage & Component & Time (s) & Device \\
    \midrule
    \multirow{3}{*}{1} & Preprocessing     &  14 & CPU \\
                       & LID (adapt.)      & 187 & GPU \\
                       & Whisper decode    & 312 & GPU \\
    \midrule
    \multirow{2}{*}{2} & IPA conversion    &   2 & CPU \\
                       & Dictionary lookup &   1 & CPU \\
    \midrule
    \multirow{3}{*}{3} & ECAPA embed.      &   8 & GPU \\
                       & MMS synthesis     & 241 & GPU \\
                       & DTW prosody       &  93 & CPU \\
    \midrule
    \multirow{2}{*}{4} & LFCC-CNN train    &  38 & GPU \\
                       & FGSM attack       &  27 & GPU \\
    \midrule
    \multicolumn{2}{l}{\textbf{Total}}    & \textbf{923} & \\
    \bottomrule
  \end{tabular}
\end{table}

Table~\ref{tab:compute} breaks down the wall-clock time for each
component.  The two dominant costs are Whisper decoding (312\,s)
and MMS synthesis (241\,s), both of which scale linearly with audio
duration.  The LID domain adaptation (187\,s) is a one-time cost;
at inference time the forward pass for the 10-minute clip takes
$<3$\,s.  DTW prosody transfer (93\,s) is the most CPU-intensive
step: its quadratic memory footprint required the 5$\times$
downsampling heuristic described in Section~\ref{sec:dtw}.  The
entire pipeline processes 600\,s of audio in approximately 923\,s
(1.54$\times$ real-time) on this hardware configuration.

%% ──────────────────────────────────────────────────────────────────
\section{Ablation Experiments}
\label{sec:ablation}

\paragraph{Decoder selection rate.}
The quality-score selector preferred the constrained decode for
approximately 61\% of utterances.  On spans containing the terms
\emph{cepstrum}, \emph{mel-filterbank}, or \emph{stochastic}, the
constrained output was selected in all cases.  For purely Hindi spans
the baseline was selected significantly more often, consistent with
the LM being English-vocabulary dominated.

\paragraph{DTW prosody transfer.}
As shown in Figure~\ref{fig:prosody}, DTW warping lifts F0 correlation
from 0.139 to 0.242 and energy correlation from 0.786 to 0.895.
The pitch block size of 120\,ms was chosen from a small grid search
over $\{50, 100, 120, 200\}$\,ms; blocks shorter than 100\,ms
introduced click artifacts from PSOLA segment boundaries, while blocks
longer than 150\,ms were perceptually discontinuous.

\paragraph{LID smoothing.}
A 3-frame median produces 31 predicted switches on the held-out set
(vs.\ 7 for 15-frame).  However, it also introduces spurious
rapid-fire language flips that cause the downstream ASR language hint
to oscillate, which doubles WER on mixed spans.  The 15-frame choice
is therefore the right engineering compromise for this pipeline's
intended use.

%% ──────────────────────────────────────────────────────────────────
\section{Limitations}
\label{sec:limits}

Three fundamental limitations are worth stating explicitly.

\paragraph{Switch-boundary recall.}
The smoothing-based LID predictor is poorly suited for boundary
localization.  A boundary-aware recurrent model trained with a
sequence-segmentation objective (e.g., connectionist temporal
classification over the switch sequence) would be a principled
replacement, but training such a model requires dense human-annotated
switch timestamps rather than the span-level annotation we have.

\paragraph{Annotation quality.}
Our evaluation rests on semi-automatic pseudo-gold annotations.
The absolute numbers are reliable enough to confirm that the system
is working as intended, but they should not be compared against
published baselines on controlled benchmarks.

\paragraph{Voice identity in TTS.}
We chose MMS Maithili over SpeechT5 for language correctness, but
this means the output voice does not resemble the student's reference
recording.  Combining a multilingual voice-cloning model with native
Maithili phonology remains an open research problem.

%% ──────────────────────────────────────────────────────────────────
\section{Conclusion}
\label{sec:conclusion}

We built and evaluated a complete code-switched speech pipeline spanning
language identification, constrained ASR, phonetic conversion, Maithili
translation, voice cloning with prosody transfer, and adversarial
robustness analysis.  The final system achieves a LID macro-F1 of
\LIDMacroFOne, English WER of \EnglishWER, Hindi WER of \HindiWER,
MCD of \MCDLectureMMS, spoofing EER of \SpoofEER, and a successful
FGSM attack at $\epsilon=\FGSMEpsilon$ with SNR\,=\,\FGSMSNR\,dB.

The principal lesson is that pipeline coupling matters.  Every
engineering decision in an early stage propagates to downstream
evaluation: smoothing in LID affects WER; chunk duration in the TTS
manifest affects the DTW alignment and therefore the MCD; and the
exact log-Mel front-end used during training must be replicated
inside the attack loop for the adversarial gradient to transfer.
Getting these couplings right was the most demanding part of the
project.

Future work should address boundary localization with a dedicated
sequence-segmentation loss and explore multilingual voice-conditioned
synthesizers that can produce native-phonology Maithili while
preserving a target speaker's voice identity.

A second promising direction is richer domain adaptation for Whisper.
The current n-gram logit bias is entirely lexical: it boosts tokens
that appear in the domain vocabulary but does not model
co-occurrence patterns.  Replacing the 3-gram JSON model with a
small neural language model fine-tuned on Speech Understanding
course transcripts would provide both bigram and trigram context,
likely closing the gap between constrained and unconstrained
decoding for Hindi-dominant spans.

Finally, the DTW prosody module currently treats pitch, energy, and
timing as three independent one-dimensional signals.  A joint
DTW over the three-channel feature matrix
(Eq.~\ref{eq:prosodyfeat}) would preserve the natural co-variation
between intonation and emphasis, potentially improving both the
F0 and energy correlations simultaneously.  Alternatively,
a learned prosody predictor trained on the (source, reference) pairs
could replace the deterministic DTW entirely and generalize to
new speakers without re-running the alignment.

\noindent\textbf{Acknowledgements.}
The teaching team of the Speech Understanding course approved
the use of the Lex Fridman - Narendra Modi podcast in place of
the originally specified lecture audio, and we thank them for
that accommodation.

%% ──────────────────────────────────────────────────────────────────
{\small
\bibliographystyle{ieeenat_fullname}
\bibliography{main}
}

\end{document}